\documentclass[11pt]{article}

\usepackage{amsmath}
\usepackage{amssymb}
\usepackage{booktabs}
\usepackage{graphicx}
\usepackage[hidelinks]{hyperref}
\usepackage[nameinlink,noabbrev]{cleveref}

\title{Transcender: Cross-Layer Composition Control for Adaptive Transformer Inference}
\author{Jun Son}
\date{March 22, 2026}

\begin{document}

\maketitle

\begin{abstract}
Transcender is best understood not as a generic early-exit method, but as a cross-layer composition control framework for adaptive transformer inference. The central claim of this repository is that adaptive-depth inference is constrained less by exit-point selection than by whether intermediate-depth and full-depth signals can be combined without corrupting the output distribution. Across three empirical tracks, this framing produces a mixed but coherent result. On GPT-OSS 20B, Track A establishes the canonical practical frontier in the repository: \texttt{L22 top1\_agree} with entropy-gated exit reaches 0.969 exact match against full-depth greedy decoding while physically skipping about 49.5\% of layers. Track B provides a scoped negative comparison baseline: a naive Gemma 3 4B-IT to GPT-OSS 20B cross-model cascade reaches only 0.026 exact match at 0.28 TPS and 20.19 GB peak memory on the local Apple MLX runtime. Track C now includes a denser follow-up beyond Gemma. On Gemma 3 4B-IT, compute-both agreement-aware composition reaches 0.807 exact match at both L31 and L20 versus 0.198 and 0.026 for fixed exit at those layers, while real selective-depth at both checkpoints remains practically weak. Focused late-checkpoint follow-up on Llama 3.1 8B and Mistral 7B reproduces the same benchmark-side pattern: fixed exit remains materially bad, compute-both \texttt{top1\_agree} recovers full-depth output, and real selective-depth still trails the matched full-depth baseline after replay/cache repair. The resulting picture is architecture-dependent: strong practical adaptive depth on the tested MoE model, but only partial methodology transfer and negative dense speed-validation across the tested dense families and local runtime.
\end{abstract}

\section{Introduction}

Variable-compute inference for transformer language models is commonly framed in one of two ways: early exit within a model, or draft-verifier decoding across models \cite{elbayad2020depthadaptive,leviathan2023speculative}. In both cases, the usual question is when the model has done enough work to emit the next token. That framing is incomplete.

This repository argues for a different abstraction: the central problem is cross-layer composition control. If an intermediate layer and the final layer disagree, a system needs a principled rule for deciding whether to trust the shallow path, the deep path, or a combination of the two. Without that rule, adaptive-depth inference either collapses into low-quality fixed exit or pays the cost of computing both paths without recovering a practical speed-quality tradeoff.

The evidence in this repository supports that reframing across three tracks. Track A establishes a practical same-model adaptive-depth frontier on GPT-OSS 20B. Track B shows that a naive cross-model cascade can fail badly when tokenizer, distribution, and orchestration costs all work against it. Track C shows both the upside and the limit of the composition-control view on dense models in this local Apple MLX runtime: compute-both agreement-aware composition recovers substantial quality on Gemma 3 4B-IT, Llama 3.1 8B, and Mistral 7B late-checkpoint follow-ups, but real selective-depth implementations still do not beat their matched full-depth baselines.

\section{Distinctive Contributions}

\begin{enumerate}
\item \textbf{The Subspace Paradox.} Hidden states from different transformer layers occupy geometrically separated subspaces in the tested models. This makes hidden-state interpolation structurally invalid for cross-layer composition.
\item \textbf{Agreement-aware logit-space composition.} Logit blending with an explicit agreement gate improves output quality across both the tested MoE and dense architectures.
\item \textbf{Empirical characterization of the cascade tax.} The repository measures how a naive cross-model cascade can degrade when tokenizer mismatch, distribution mismatch, and orchestration overhead compound.
\item \textbf{Family-sensitive dense depth profiles.} Gemma 3 4B-IT exhibits a late plateau, while focused recon on Llama 3.1 8B and Mistral 7B shows different plateau zones and resolution entry points. The benchmark-side limitation generalizes more cleanly than Gemma's exact internal profile shape.
\item \textbf{Real dense selective-depth speed validation.} Continuation-safe selective-depth runtimes were implemented at L31 and L20 on Gemma 3 4B-IT and at late checkpoints on Llama 3.1 8B and Mistral 7B. All physically skipped deep layers on some decode tokens, but none recovered a practical speed-quality frontier on this local Apple MLX runtime.
\end{enumerate}

\section{Method}

\subsection{Cross-Layer Geometric Separation}

Let $h_i$ and $h_j$ denote hidden states from layers $i$ and $j$ with $i < j$. When projected into a shared PCA space, the repository measures cross-layer separation using
\begin{equation}
\mathrm{SeparationRatio}(i,j) =
\frac{\lVert c_i - c_j \rVert_2}{\max(\sigma_i,\sigma_j)},
\label{eq:separation}
\end{equation}
where $c_i$ and $c_j$ are cluster centroids and $\sigma_i$, $\sigma_j$ are within-cluster spreads. Ratios of 4.11 on GPT-OSS 20B and 2.46 on Gemma 3 4B-IT indicate that hidden states from different layers are not safely interpolable in the tested settings.

\subsection{Agreement-Aware Logit Composition}

Instead of blending hidden states, Transcender projects intermediate and final states into a shared logit space and composes them there. Let $\ell_t^e$ be logits at exit layer $e$ for token step $t$, and let $\ell_t^N$ be logits at full depth $N$. The composition rule is
\begin{equation}
\ell_t^{\mathrm{out}} =
\begin{cases}
(1-\alpha)\,\ell_t^N + \alpha\,\ell_t^e, &
\arg\max \ell_t^e = \arg\max \ell_t^N, \\
\ell_t^N, & \text{otherwise.}
\end{cases}
\label{eq:blend}
\end{equation}
This is the \texttt{top1\_agree} rule used throughout the repository. Its purpose is conservative: blend only when shallow and deep paths already agree on token identity, otherwise discard the shallow logits.

\subsection{Selective-Depth Runtime}

The dense selective-depth runtime used in Track C differs from the compute-both composition path. Prompt prefill is always run at full depth so the deep cache is valid at decode start. During autoregressive decode, layers $0 \ldots e$ are run first, a confidence score $s(\ell_t^e)$ is computed, and an early token is accepted only when confidence exceeds a threshold:
\begin{equation}
a_t = \mathbf{1}[s(\ell_t^e) \ge \tau], \qquad
\hat{y}_t =
\begin{cases}
\arg\max \ell_t^e, & a_t = 1, \\
\arg\max \ell_t^N, & a_t = 0.
\end{cases}
\label{eq:selective}
\end{equation}
If later continuation is required after one or more accepted shallow tokens, the previously skipped decode tokens are replayed through the deep layers to restore deep-cache correctness. In the Track C extension, this replay cost is part of the measured runtime and is essential to the interpretation of the result.

\section{Experimental Setup}

All results were collected on a single Apple M1 Pro system with 32 GB unified memory using the Apple MLX runtime. The prompt family is fixed across the repository: five factual or explanatory prompts on quantum entanglement, the French Revolution, recursion, TCP versus UDP, and photosynthesis. Each prompt generates 48 tokens with greedy decoding. Aggregate metrics exclude the first prompt as warmup in order to reduce distortion from JIT compilation and initial memory allocation.

The primary quality metric is exact token-sequence match against a full-depth baseline on the same model. Secondary metrics include generation throughput in tokens per second, time to first token, elapsed time, peak memory, and layer-saving statistics where applicable.

The evaluated models are GPT-OSS 20B, a 24-layer Mixture-of-Experts model run in MXFP4 quantization; Gemma 3 4B-IT, a 34-layer dense model in the default MLX configuration \cite{gemma2024report}; and MLX 4bit dense follow-up models from the Llama and Mistral families: Llama 3.1 8B Instruct and Mistral 7B Instruct v0.3. Track B combines Gemma 3 4B-IT and GPT-OSS 20B in a naive cross-model cascade. Track C evaluates Gemma 3 4B-IT in both the original six-mode adaptive benchmark and dedicated selective-depth runtimes, then adds focused late-checkpoint validation on Llama and Mistral.

\section{Results}

\subsection{Track A: Canonical Same-Model Adaptive Depth}

\begin{table}[t]
\centering
\small
\caption{Track A on GPT-OSS 20B.}
\label{tab:tracka}
\begin{tabular}{@{}lccc@{}}
\toprule
Configuration & Gen TPS & Exact Match & Avg. Layers Saved \\
\midrule
Full depth (L23) & 31.79 & 1.000 & 0\% \\
L22 \texttt{top1\_agree} & 27.41 & 0.969 & 49.5\% \\
L23 \texttt{top1\_agree} & 32.41 & 1.000 & 0\% \\
\bottomrule
\end{tabular}
\end{table}

Track A is the canonical practical result in the repository. As shown in \cref{tab:tracka}, \texttt{L22 top1\_agree} preserves 0.969 exact match while physically skipping about half of the model's layers on average. This is not a compute-both artifact. The \texttt{avg\_layers\_saved} value reflects real entropy-gated skipping in the MLX runtime.

\subsection{Track B: Scoped Negative Cascade Baseline}

\begin{table}[t]
\centering
\small
\caption{Track B comparison on the local Apple MLX runtime. The Track A row is included only as the comparison point used in this repository.}
\label{tab:trackb}
\begin{tabular}{@{}lcccc@{}}
\toprule
Mode & Gen TPS & Exact Match & Peak Mem (GB) & TTFT (s) \\
\midrule
Draft only (Gemma 3 4B-IT) & 18.99 & 0.006 & 7.28 & 0.23 \\
GPT-OSS full depth & 30.94 & 1.000 & 12.96 & 0.66 \\
Naive cascade & 0.28 & 0.026 & 20.19 & 11.83 \\
Track A L22 \texttt{top1\_agree} & 20.22 & 0.969 & 12.96 & 0.98 \\
\bottomrule
\end{tabular}
\end{table}

Track B is not a claim about cascade systems in general. It is a negative comparison baseline for this naive implementation, this Gemma 3 4B-IT to GPT-OSS 20B model pair, and this local Apple MLX runtime. In \cref{tab:trackb}, the naive cascade reaches 0.026 exact match at 0.28 TPS and 20.19 GB peak memory. The likely causes are tokenizer mismatch, distribution mismatch, and two-model orchestration overhead. The repository therefore treats Track B as evidence about the cascade tax in this specific setup, not as a universal statement against speculative decoding \cite{leviathan2023speculative}.

\subsection{Track C: Fixed Exit Failure and Compute-Both Composition Signal}

\begin{table}[t]
\centering
\small
\caption{Original Track C adaptive benchmark on Gemma 3 4B-IT.}
\label{tab:trackc_adaptive}
\begin{tabular}{@{}lcccc@{}}
\toprule
Mode & Gen TPS & Exact Match & Layers Saved & Agreement Rate \\
\midrule
Full depth L33 & 15.01 & 1.000 & 0\% & --- \\
Fixed exit L16 & 22.81 & 0.010 & 50\% & --- \\
Fixed exit L20 & 20.81 & 0.026 & 38\% & --- \\
Fixed exit L31 & 14.95 & 0.198 & 6\% & --- \\
\texttt{top1\_agree} L31 & 7.19 & 0.807 & 0\% & 81.9\% \\
Naive blend L31 & 10.01 & 0.688 & 0\% & 100\% \\
\bottomrule
\end{tabular}
\end{table}

Track C first established that fixed exit is catastrophic on the tested dense model. Even the latest tested fixed exit, L31, saves only 6\% of layers and falls to 0.198 exact match. The same benchmark also established a genuine composition signal: compute-both agreement-aware composition at L31 reaches 0.807 exact match, well above both fixed exit and naive blending in \cref{tab:trackc_adaptive}. This is the core dense-model evidence supporting the composition-control thesis.

\subsection{Track C: Real Selective-Depth Speed Validation at L31 and L20}

\begin{table}[t]
\centering
\small
\caption{Dedicated Track C selective-depth benchmark. These numbers should be interpreted relative to the full-depth row in the same table because they come from a separate selective-depth runtime.}
\label{tab:trackc_selective}
\begin{tabular}{@{}p{3.4cm}ccccc@{}}
\toprule
Mode & Gen TPS & Exact Match & Accept & Skip & Avg. Layers Saved \\
\midrule
Full depth L33 & 20.05 & 1.000 & --- & 0\% & 0\% \\
Fixed exit L31 & 19.10 & 0.198 & --- & --- & 5.9\% \\
\texttt{top1\_agree} compute-both L31 & 16.31 & 0.807 & 81.9\% & 0\% & 0\% \\
Selective margin L31 & 18.26 & 0.229 & 78.2\% & 3.7\% & 0.2\% \\
Selective entropy L31 & 18.20 & 0.208 & 93.6\% & 34.0\% & 2.0\% \\
\bottomrule
\end{tabular}
\end{table}

The selective-depth extension answers a narrower question than \cref{tab:trackc_adaptive}: can the late-layer agreement signal at L31 be converted into real skipping? The answer is yes in an operational sense and no in a practical sense. As shown in \cref{tab:trackc_selective}, the runtime physically skipped layers 32 and 33 on some decode tokens. The entropy rule achieved a 34.0\% realized skip rate over continuation-eligible decode steps. However, only two of thirty-four layers were available to skip, so average layer savings remained just 2.0\% in the best case. Neither selective mode beat the dedicated full-depth baseline of 20.05 TPS, and both remained close to fixed-exit quality.

This distinction matters. Compute-both \texttt{top1\_agree} is a quality-control result with no physical skip. Real selective depth is a speed-validation result. On Gemma 3 4B-IT at L31, with this implementation and this local Apple MLX runtime, replay and control overhead erased the small available compute savings.

\begin{table}[t]
\centering
\small
\caption{Dedicated L20 selective-depth probe.}
\label{tab:trackc_selective_l20}
\begin{tabular}{@{}p{3.4cm}ccccc@{}}
\toprule
Mode & Gen TPS & Exact Match & Accept & Skip & Avg. Layers Saved \\
\midrule
Full depth L33 & 15.16 & 1.000 & --- & 0\% & 0\% \\
Fixed exit L20 & 20.20 & 0.026 & --- & --- & 38.2\% \\
\texttt{top1\_agree} compute-both L20 & 10.54 & 0.807 & 28.2\% & 0\% & 0\% \\
Selective margin L20 & 11.29 & 0.083 & 70.2\% & 13.8\% & 5.3\% \\
Selective entropy L20 & 13.26 & 0.073 & 73.4\% & 41.0\% & 15.7\% \\
Selective hybrid L20 & 11.37 & 0.125 & 61.2\% & 6.4\% & 2.4\% \\
\bottomrule
\end{tabular}
\end{table}

The L20 probe answers the obvious next question after L31: does a materially larger skip budget recover a practical frontier? \Cref{tab:trackc_selective_l20} shows that it does not in this configuration. The L20 entropy rule increased average layers saved to 15.7\% and the realized skip rate to 41.0\%, but still reached only 13.26 TPS against a matched full-depth baseline of 15.16 TPS. Quality also collapsed toward fixed-exit behavior at 0.073 exact match. Margin and hybrid variants were similarly weak. L20 therefore strengthens the failure boundary rather than overturning it: larger skip budget alone does not solve dense selective-depth when continuation quality is poor and replay cost remains significant.

\begin{table}[t]
\centering
\small
\caption{Focused late-checkpoint dense-family follow-up on the local Apple MLX runtime.}
\label{tab:trackc_multifamily}
\begin{tabular}{@{}p{3.0cm}cccccc@{}}
\toprule
Model & Late CP & Fixed Exit & Compute-Both & Selective & Selective TPS & Avg. Saved \\
\midrule
Llama 3.1 8B & L29 & 0.151 & 1.000 & 0.469 & 16.60 vs 17.81 & 0.3\% \\
Mistral 7B v0.3 & L29 & 0.109 & 1.000 & 0.318 & 17.25 vs 29.81 & 0.2\% \\
\bottomrule
\end{tabular}
\end{table}

\subsection{Track C: Focused Multi-Family Dense Follow-Up}

\Cref{tab:trackc_multifamily} summarizes the late-checkpoint follow-up on Llama 3.1 8B and Mistral 7B. The benchmark-side pattern from Gemma reproduces cleanly: late fixed exit fails materially, compute-both \texttt{top1\_agree} recovers full-depth output, and real selective-depth physically skips layers but still does not recover a practical speed-quality frontier. On Llama, fixed exit at L29 reached only 0.151 exact match, while compute-both \texttt{top1\_agree} reached 1.000 and real selective-depth entropy reached 0.469 exact at 16.60 TPS versus a matched full-depth baseline of 17.81 TPS. On Mistral, fixed exit at L29 reached 0.109 exact match, while compute-both \texttt{top1\_agree} again reached 1.000 and real selective-depth entropy reached 0.318 exact at 17.25 TPS versus 29.81 TPS full depth.

These follow-up runs materially broaden the dense benchmark result beyond Gemma alone, but they do not justify calling Gemma's exact internal profile universal. The focused recon runs were family-sensitive: Gemma showed a late L20--L29 plateau, Llama showed a milder middle plateau around L12--L14 with heuristic resolution entry at L27, and Mistral showed an earlier plateau around L2--L8 with heuristic resolution entry at L26. The benchmark limitation therefore generalizes more cleanly than the exact KL-profile geometry.

\section{Limitations}

\begin{itemize}
\item \textbf{Architecture sensitivity.} The strongest practical result appears on the tested MoE model, not on the tested dense model. The magnitude of adaptive-depth gains is architecture-dependent in this repository.
\item \textbf{Track B scope.} The cascade result is specific to a naive implementation, the Gemma 3 4B-IT and GPT-OSS 20B model pair, and the local Apple MLX runtime. It should not be generalized to speculative decoding as a whole.
\item \textbf{Local runtime dependence.} All throughput, latency, and memory numbers are specific to a single Apple M1 Pro system using MLX. They should not be compared directly to GPU-based inference reports.
\item \textbf{Small prompt suite.} The evaluation uses five prompts with warmup exclusion. That is adequate for directional findings, but not for statistical confidence.
\item \textbf{Dense selective-depth remains unsuccessful here.} Real selective-depth runtimes on Gemma, Llama, and Mistral were operationally correct, but none recovered a practical speed-quality frontier in this configuration.
\item \textbf{Limited continuation mechanisms.} The completed dense selective-depth results show that moving the checkpoint earlier or relying on late-checkpoint confidence alone is not enough by itself. Better continuation criteria, token-difficulty-aware routing, and lower-overhead cache repair remain open engineering and research questions.
\item \textbf{Family-sensitive internal structure.} The benchmark-side dense limitation generalized across the tested families, but the internal KL-profile shape did not. The paper therefore does not claim a single universal dense depth profile.
\end{itemize}

\section{Conclusion}

The evidence in this repository supports a narrow but consistent conclusion: Transcender is best understood not as a generic early-exit method, but as a cross-layer composition control framework for adaptive transformer inference.

Track A remains the canonical result. On GPT-OSS 20B, \texttt{L22 top1\_agree} with entropy-gated exit achieves 0.969 exact match with about 49.5\% real layer savings. Track B remains a scoped negative comparison baseline for this naive cascade implementation, this Gemma 3 4B-IT to GPT-OSS 20B model pair, and this local Apple MLX runtime. Track C now has three layers: fixed-exit failure, compute-both composition signal, and weak real selective-depth. Gemma provides the deepest dense case study, and focused late-checkpoint follow-up on Llama 3.1 8B and Mistral 7B reproduces the same benchmark-side pattern. The dense-model conclusion is therefore not that adaptive depth is impossible, but that neither fixed exit nor the current real selective-depth mechanisms establish a practical frontier in the tested dense follow-ups.

The next research question is correspondingly narrower: whether selective-depth strategies with better, potentially family-sensitive continuation criteria, token-difficulty-aware routing, and lower-overhead cache-aware continuation can recover a practical frontier on dense models. That remains to be validated.

\appendix

\section{Reproduction Commands}
\label{app:repro}

\noindent All commands assume execution from the repository root.

\begin{itemize}
\item \texttt{python scripts/track\_a/transcender\_exit\_layer\_benchmark.py}
\item \texttt{python scripts/track\_b/transcender\_track\_b\_benchmark.py --large-model /path/to/gpt-oss-20b-raw --draft-model /path/to/gemma-3-4b-it}
\item \texttt{python scripts/track\_c/transcender\_track\_c\_gemma\_profile.py --model /path/to/gemma-3-4b-it}
\item \texttt{python scripts/track\_c/transcender\_track\_c\_gemma\_benchmark.py --model /path/to/gemma-3-4b-it --early-layer 16 --mid-layer 20 --late-layer 31}
\item \texttt{python scripts/track\_c/transcender\_track\_c\_gemma\_selective\_depth.py --model /path/to/gemma-3-4b-it --late-layer 31}
\item \texttt{python scripts/track\_c/transcender\_track\_c\_gemma\_selective\_depth.py --model /path/to/gemma-3-4b-it --late-layer 20 --margin-threshold 0.02 --entropy-threshold 0.65 --include-hybrid-mode}
\item \texttt{python scripts/track\_c/transcender\_track\_c\_dense\_family\_validation.py --model /path/to/llama-3.1-8b-instruct-4bit --family llama --profile-output artifacts/recon/recon\_llama3\_8b.json --benchmark-output artifacts/dense\_followup/transcender\_track\_c\_llama3\_8b\_results.json --late-layer 29 --entropy-threshold 0.15}
\item \texttt{python scripts/track\_c/transcender\_track\_c\_dense\_family\_validation.py --model /path/to/mistral-7b-instruct-v0.3-4bit --family mistral --profile-output artifacts/recon/recon\_mistral7b.json --benchmark-output artifacts/dense\_followup/transcender\_track\_c\_mistral7b\_results.json --late-layer 29 --entropy-threshold 0.15}
\end{itemize}

\section{Benchmark Artifacts}
\label{app:artifacts}

\begin{itemize}
\item \texttt{artifacts/track\_a/transcender\_exit\_layer\_benchmark.json}: Track A per-prompt and aggregate metrics for L22 and L23 configurations.
\item \texttt{artifacts/track\_b/transcender\_track\_b\_benchmark.json}: Track B per-prompt and aggregate metrics for draft-only, full-depth, cascade, and Track A comparison modes.
\item \texttt{artifacts/track\_c/transcender\_track\_c\_gemma\_results.json}: Track C per-prompt and aggregate metrics for the original six-mode Gemma benchmark.
\item \texttt{artifacts/track\_c/transcender\_track\_c\_gemma\_selective\_depth\_results.json}: Track C per-prompt and aggregate metrics for real selective-depth speed validation at L31.
\item \texttt{artifacts/track\_c/transcender\_track\_c\_gemma\_selective\_depth\_L20\_results.json}: Track C per-prompt and aggregate metrics for the dedicated L20 selective-depth probe.
\item \texttt{artifacts/recon/recon\_llama3\_8b.json}: Focused KL profile and heuristic summary for Llama 3.1 8B Instruct 4bit.
\item \texttt{artifacts/dense\_followup/transcender\_track\_c\_llama3\_8b\_results.json}: Focused late-checkpoint validation metrics for Llama 3.1 8B Instruct 4bit.
\item \texttt{artifacts/recon/recon\_mistral7b.json}: Focused KL profile and heuristic summary for Mistral 7B Instruct v0.3 4bit.
\item \texttt{artifacts/dense\_followup/transcender\_track\_c\_mistral7b\_results.json}: Focused late-checkpoint validation metrics for Mistral 7B Instruct v0.3 4bit.
\item \texttt{artifacts/track\_c/gemma\_kl\_profile.json}: Per-layer KL divergence values across the 34-layer Gemma 3 4B-IT stack.
\end{itemize}

\bibliographystyle{plain}
\bibliography{references}

\end{document}
